\newpage

\section{Выполнение лабораторной работы} 

\subsection*{\textbf{Часть 1. Подготовка и первичный осмотр данных}}
\addcontentsline{toc}{subsection}{Часть 1. Подготовка и первичный осмотр данных}
 
На первом этапе была выполнена загрузка датасета Wine Quality, проведён его первичный осмотр, определены размерность и типы данных, проверено наличие пропущенных значений и дубликатов. Кроме того, для каждого признака были определены тип и шкала измерений, а также проверена корректность вычисления среднего значения.
 
\subsubsection*{Загрузка данных и первичный осмотр}
 
Для работы с данными были подключены библиотеки
\texttt{pandas}, \texttt{numpy}, \texttt{matplotlib},
\texttt{seaborn} и \texttt{scipy}. Датасет был загружен из файла \texttt{wine\_quality\_merged.csv}.
 
\begin{Code}
import pandas as pd
import numpy as np
import matplotlib.pyplot as plt
import seaborn as sns
from scipy import stats
from IPython.display import display

# Загрузка данных
df = pd.read_csv('wine_quality_merged.csv')

# Первичный осмотр
print("Размерность:", df.shape)
print("Количество строк:", df.shape[0], "| Количество столбцов:", df.shape[1])
print()
print("Типы данных:")
print(df.dtypes)
print()
print('Первые пять строк:')
display(df.head())
print('Последние пять строк:')
display(df.tail())
\end{Code}

В результате выполнения кода было установлено, что размерность датасета составляет $(6497, 13)$. Это означает, что набор данных содержит 6497 наблюдений и 13 признаков.
 
Первые пять строк датасета представлены в таблице
\ref{tab:first-rows-wine}, а последние пять строк --- в таблице
\ref{tab:last-rows-wine}.
 

\begin{table}[ht]
	\centering
	\caption{Первые пять строк датасета \texttt{Wine Quality}}
	\label{tab:first-rows-wine}
	\scriptsize
	\setlength{\tabcolsep}{3pt}
	\begin{tabularx}{\textwidth}{c|CCCCCCC@{\hspace{-4pt}}C@{\hspace{-8pt}}C@{\hspace{-2pt}}CC@{\hspace{-1pt}}C@{\hspace{-6pt}}C}
		\hline
		\makecell{id} &
		\makecell{\texttt{fixed}\\\texttt{acidity}} &
		\makecell{\texttt{volatile}\\\texttt{acidity}} &
		\makecell{\texttt{citric}\\\texttt{acid}} &
		\makecell{\texttt{residual}\\\texttt{sugar}} &
		\makecell{{chlorides}} &
		\makecell{\texttt{free}\\\texttt{sulfur}\\\texttt{dioxide}} &
		\makecell{\texttt{total}\\\texttt{sulfur}\\\texttt{dioxide}} &
		\makecell{\texttt{density}} &
		\makecell{\texttt{pH}} &
		\makecell{{sulphates}} &
		\makecell{\texttt{alcohol}} &
		\makecell{\texttt{quality}} &
		\makecell{\texttt{type}} \\
		\hline
		0 & 7{,}4  & 0{,}70 & 0{,}00 & 1{,}9 & 0{,}076 & 11{,}0 & 34{,}0 & 0{,}9978 & 3{,}51 & 0{,}56 & 9{,}4 & 5 & \textit{red} \\
		1 & 7{,}8  & 0{,}88 & 0{,}00 & 2{,}6 & 0{,}098 & 25{,}0 & 67{,}0 & 0{,}9968 & 3{,}20 & 0{,}68 & 9{,}8 & 5 & \textit{red} \\
		2 & 7{,}8  & 0{,}76 & 0{,}04 & 2{,}3 & 0{,}092 & 15{,}0 & 54{,}0 & 0{,}9970 & 3{,}26 & 0{,}65 & 9{,}8 & 5 & \textit{red} \\
		3 & 11{,}2 & 0{,}28 & 0{,}56 & 1{,}9 & 0{,}075 & 17{,}0 & 60{,}0 & 0{,}9980 & 3{,}16 & 0{,}58 & 9{,}8 & 6 & \textit{red} \\
		4 & 7{,}4  & 0{,}70 & 0{,}00 & 1{,}9 & 0{,}076 & 11{,}0 & 34{,}0 & 0{,}9978 & 3{,}51 & 0{,}56 & 9{,}4 & 5 & \textit{red} \\
		\hline
	\end{tabularx}
\end{table}

\begin{table}[ht]
	\centering
	\caption{Последние пять строк датасета \texttt{Wine Quality}}
	\label{tab:last-rows-wine}
	\scriptsize
	\setlength{\tabcolsep}{3pt}
	\begin{tabularx}{\textwidth}{c|CCCC@{\hspace{1pt}}C@{\hspace{-4pt}}C@{\hspace{-1pt}}C@{\hspace{-4pt}}C@{\hspace{-9pt}}C@{\hspace{-2pt}}CC@{\hspace{-1pt}}C@{\hspace{-6pt}}C}
		\hline
		\makecell{id} &
		\makecell{\texttt{fixed}\\\texttt{acidity}} &
		\makecell{\texttt{volatile}\\\texttt{acidity}} &
		\makecell{\texttt{citric}\\\texttt{acid}} &
		\makecell{\texttt{residual}\\\texttt{sugar}} &
		\makecell{{chlorides}} &
		\makecell{\texttt{free}\\\texttt{sulfur}\\\texttt{dioxide}} &
		\makecell{\texttt{total}\\\texttt{sulfur}\\\texttt{dioxide}} &
		\makecell{\texttt{density}} &
		\makecell{\texttt{pH}} &
		\makecell{{sulphates}} &
		\makecell{\texttt{alcohol}} &
		\makecell{\texttt{quality}} &
		\makecell{\texttt{type}} \\
		\hline
		6492 & 6{,}2 & 0{,}21 & 0{,}29 & 1{,}6 & 0{,}039 & 24{,}0  & 92{,}0  & 0{,}99114 & 3{,}27 & 0{,}50 & 11{,}2 & 6 & \textit{white} \\
		6493 & 6{,}6 & 0{,}32 & 0{,}36 & 8{,}0 & 0{,}047 & 57{,}0  & 168{,}0 & 0{,}99490 & 3{,}15 & 0{,}46 & 9{,}6  & 5 & \textit{white} \\
		6494 & 6{,}5 & 0{,}24 & 0{,}19 & 1{,}2 & 0{,}041 & 30{,}0  & 111{,}0 & 0{,}99254 & 2{,}99 & 0{,}46 & 9{,}4  & 6 & \textit{white} \\
		6495 & 5{,}5 & 0{,}29 & 0{,}30 & 1{,}1 & 0{,}022 & 20{,}0  & 110{,}0 & 0{,}98869 & 3{,}34 & 0{,}38 & 12{,}8 & 7 & \textit{white} \\
		6496 & 6{,}0 & 0{,}21 & 0{,}38 & 0{,}8 & 0{,}020 & 22{,}0  & 98{,}0  & 0{,}98941 & 3{,}26 & 0{,}32 & 11{,}8 & 6 & \textit{white} \\
		\hline
	\end{tabularx}
\end{table}

\subsubsection*{Проверка пропущенных значений и дубликатов}
 
Для проверки качества данных был использован следующий код:
 
\begin{Code}
print("Пропуски по столбцам:")
print(df.isnull().sum())
print()

print("Всего пропусков:", df.isnull().sum().sum())
print()

print("Количество дубликатов:", df.duplicated().sum())
print()

print("Дубликаты:")
display(df[df.duplicated()])
\end{Code}
 
В результате проверки было установлено, что пропущенные значения
отсутствуют во всех столбцах датасета. Общее количество пропусков
равно нулю. При этом было обнаружено одно дублирующееся наблюдение.
 
\begin{table}[ht]
    \centering
    \caption{Результаты проверки пропусков и дубликатов}
    \label{tab:missing-duplicates-wine}
    \begin{tabularx}{0.65\textwidth}{LC}
        \hline
        Показатель & Значение \\
        \hline
        Общее количество пропусков & 0 \\
        Количество дубликатов & 1177 \\
        \hline
    \end{tabularx}
\end{table}

В результате проверки было установлено, что пропущенные значения отсутствуют во всех столбцах датасета. Общее количество пропусков равно нулю. При этом было обнаружено 1177 дублирующихся наблюдений. На данном этапе
дубликаты были зафиксированы, но не удалялись, поскольку целью работы
является первичный разведочный анализ исходного датасета.
 
\subsubsection*{Определение типов признаков}
 
Для определения типов данных и количества уникальных значений в каждом
столбце использовался следующий код:
 
\begin{Code}
for col in df.columns:
	print(f"{col}: тип = {df[col].dtype}, "
		f"уникальных значений = {df[col].nunique()}")

print()
print("Классы в 'type':")
print(df['type'].value_counts())
\end{Code}
 
Результаты выполнения кода представлены в таблицах
\ref{tab:feature-types-wine} и \ref{tab:species-count-wine}.
 
\begin{table}[ht]
    \centering
    \caption{Типы и количество уникальных значений признаков}
    \label{tab:feature-types-wine}
    \begin{tabularx}{\textwidth}{LCC}
        \hline
        Признак & Тип данных & Количество уникальных значений \\
        \hline
        \texttt{fixed acidity} & \texttt{float64} & 106 \\
        \texttt{volatile acidity}  & \texttt{float64} & 187 \\
        \texttt{citric acid} & \texttt{float64} & 89 \\
        \texttt{residual sugar}  & \texttt{float64} & 316 \\
        \texttt{chlorides}  & \texttt{float64}     & 214  \\
        \texttt{free sulfur dioxide}       & \texttt{float64}     & 135  \\
        \texttt{total sulfur dioxide}       & \texttt{float64}     & 276  \\
        \texttt{density}       & \texttt{float64}     & 998  \\
        \texttt{pH}       & \texttt{float64}     & 108  \\
        \texttt{sulphates}       & \texttt{float64}     & 111  \\
        \texttt{alcohol}       & \texttt{float64}     & 111  \\
        \texttt{quality}       & \texttt{int64}     & 7  \\
        \texttt{type}       & \texttt{str}     & 2  \\
        \hline
    \end{tabularx}
\end{table}
 
\begin{table}[ht]
    \centering
    \caption{Распределение наблюдений по типам вина}
    \label{tab:species-count-wine}
    \begin{tabularx}{0.65\textwidth}{LC}
        \hline
        Классы в 'type': & Количество наблюдений \\
        \hline
        \textit{white} & 4898 \\
        \textit{red} & 1599 \\
        \hline
    \end{tabularx}
\end{table}
\newpage

Таким образом, датасет содержит 11 количественных признаков и 2 категориальных признака (\texttt{type} и \texttt{quality}). Признак \texttt{quality} является порядковой переменной, так как оценки качества имеют естественный порядок. Признак \texttt{type} — номинальный. Наблюдения распределены между двумя типами вина: красное и белое.
 
\subsubsection*{Типы и шкалы измерений признаков}
 
Количественные признаки (например, \texttt{fixed acidity}, \texttt{volatile acidity}, \texttt{alcohol}) представляют собой измерения физико-химических свойств и относятся к шкале отношений. Признак \texttt{quality} является категориальным порядковым признаком. Признак type является категориальным номинальным признаком.

\begin{table}[ht]
    \centering
    \caption{Типы и шкалы измерений признаков}
    \label{tab:scales-wine}
    \scriptsize
    \begin{tabularx}{\textwidth}{LLL}
        \hline
        Признак & Тип признака & Шкала измерений \\
        \hline
        \texttt{fixed acidity} &
        количественный непрерывный &
        отношений \\
 
        \texttt{volatile acidity} &
        количественный непрерывный &
        отношений \\
 
        \texttt{citric acid} &
        количественный непрерывный &
        отношений \\
 
        \texttt{residual sugar} &
        количественный непрерывный &
        отношений \\
        
        \texttt{chlorides} &
        количественный непрерывный &
        отношений \\
        
        \texttt{free sulfur dioxide} &
        количественный непрерывный &
        отношений \\
        
        \texttt{total sulfur dioxide} &
        количественный непрерывный &
        отношений \\
        
        \texttt{density} &
        количественный непрерывный &
        отношений \\
        
        \texttt{pH} &
        количественный непрерывный &
        отношений \\
        
         \texttt{sulphates} &
        количественный непрерывный &
        отношений \\
        
         \texttt{alcohol} &
        количественный непрерывный &
        отношений \\
        
        \texttt{quality} &
        количественный дискретный &
        порядковая \\
 
        \texttt{type} &
        категориальный номинальный &
        номинальная \\
        \hline
    \end{tabularx}
\end{table}
 
\subsubsection*{Проверка корректности вычисления среднего}
 
Для проверки центральной тенденции использовалась переменная \texttt{features}, содержащая названия количественных признаков:
 
\begin{Code}
# Анализ корректности вычисления среднего для каждого признака
analysis = pd.DataFrame({
	'Среднее': df[features].mean(),
	'Медиана': df[features].median(),
	'Мода': df[features].mode().iloc[0],
	'Ст. откл.': df[features].std(ddof=1),
	'Асимметрия': df[features].skew(),
	'Эксцесс': df[features].kurt(),
})

# абсолютное и относительное расхождение среднего и медианы
analysis['|Среднее − Медиана|'] = (analysis['Среднее'] - analysis['Медиана']).abs()
analysis['Расхождение, %'] = (
analysis['|Среднее − Медиана|'] / analysis['Медиана'].abs() * 100
)

# коэффициент вариации
analysis['CV, %'] = analysis['Ст. откл.'] / analysis['Среднее'] * 100

# оценка корректности среднего
def judge(row):
if row['Расхождение, %'] < 5 and abs(row['Асимметрия']) < 0.5:
return 'Среднее корректно (симметричное распределение)'
elif row['Расхождение, %'] < 15 and abs(row['Асимметрия']) < 1:
return 'Среднее приемлемо, есть умеренная асимметрия'
else:
return 'Среднее смещено (асимметрия / влияние выбросов)'

analysis['Вывод'] = analysis.apply(judge, axis=1)

display(analysis.round(3))
\end{Code}
 
Среднее арифметическое характеризует общий уровень значений признака. Медиана делит упорядоченный ряд наблюдений на две равные части, а мода показывает наиболее часто встречающееся значение.
 
\begin{table}[ht]
    \centering
    \caption{Средние значения количественных признаков}
    \label{tab:means-wine}
    \begin{tabularx}{0.65\textwidth}{LCCC}
        \hline
        Признак & Среднее значение & Медиана & Мода \\
        \hline
        \texttt{fixed acidity} & 7{,}215 & 7{,}000 & 6{,}800\\
        \texttt{volatile acidity}  & 0{,}340 & 0{,}290 & 0{,}280\\
        \texttt{citric acid} & 0{,}319 & 0{,}310 & 0{,}300\\
        \texttt{residual sugar}  & 5{,}443 & 3{,}000 & 2{,}000\\
        \texttt{chlorides}  & 0{,}056 & 0{,}047 & 0{,}044\\
        \texttt{free sulfur dioxide}  & 30{,}525 & 29{,}000 & 29{,}000\\
        \texttt{total sulfur dioxide}  & 115{,}745 & 118{,}000 & 111{,}000\\
        \texttt{desity}  & 0{,}995 & 0{,}995 & 0{,}997\\
        \texttt{pH}  & 3{,}219 & 3{,}210 & 3{,}160\\
        \texttt{sulphates}  & 0{,}531 & 0{,}510 & 0{,}500\\
        \texttt{alcohol}  & 10{,}492 & 10{,}300 & 9{,}500\\
        \hline
    \end{tabularx}
\end{table}
 
\subsubsection*{Вывод по части 1}
 
В результате первичного осмотра было установлено, что датасет Wine Quality содержит 6497 наблюдений и 13 признаков. Одиннадцать признаков являются количественными непрерывными признаками, измеренными по шкале отношений, признак \texttt{quality} --- количественный дискретный, а признак \texttt{type} является номинальным категориальным признаком.

Пропущенные значения в датасете отсутствуют, однако обнаружено 1177 дублирующихся наблюдений. Вычисление среднего значения корректно только для количественных признаков. Полученные данные пригодны для дальнейшего разведочного анализа и построения статистических визуализаций.
 
\subsection*{\textbf{Часть 2. Проверка базовых предположений (4-Plot)}}
\addcontentsline{toc}{subsection}{Часть 2. Проверка базовых предположений (4-Plot)}

В данной части была выполнена проверка четырёх фундаментальных предположений измерительного процесса с помощью графиков последовательности наблюдений, лаговых диаграмм, гистограмм и Q–Q-графиков.


\subsubsection*{Run Sequence Plot}

Run Sequence Plot используется для анализа изменения значений признака в порядке следования наблюдений, а также для проверки постоянства среднего уровня и вариации.

\begin{Code}
fig, axes = plt.subplots(4, 3, figsize=(15, 14))
axes = axes.flatten()
for i, col in enumerate(features):
	ax = axes[i]
	ax.plot(df[col], marker='.', markersize=3, linewidth=0.8)
	ax.set_title(f'Run Sequence Plot: {col}')
	ax.set_xlabel('Индекс наблюдения')
	ax.set_ylabel('Значение')
plt.tight_layout()
plt.show()
\end{Code}

\begin{figure}[H]
    \centering
    \includegraphics[width=0.9\textwidth]{run_sequence_plot.png}
    \caption{Графики последовательности наблюдений}
    \label{fig:run-sequence}
\end{figure}

На графиках последовательности наблюдений видны изменения среднего уровня. Это связано с тем, что в датасете сначала идут наблюдения для одного типа вина, а затем для другого. Следовательно, общий средний уровень признаков не является постоянным для всей выборки. Предположение о фиксированном положении нарушается.

\subsubsection*{Lag Plot}

Lag Plot используется для анализа зависимости между соседними
наблюдениями. В данном случае сравниваются значения признака в текущем
наблюдении $x[t]$ и следующем наблюдении $x[t+1]$.

\begin{Code}
fig, axes = plt.subplots(4, 3, figsize=(15, 14))
axes = axes.flatten()
for i, col in enumerate(features):
	ax = axes[i]
	ax.scatter(df[col][:-1], df[col][1:], s=10, alpha=0.6)
	ax.set_title(f'Lag Plot: {col}')
	ax.set_xlabel('x[t]')
	ax.set_ylabel('x[t+1]')
	ax.set_aspect('equal')
plt.tight_layout()
plt.show()
\end{Code}

\begin{figure}[H]
    \centering
    \includegraphics[width=0.9\textwidth]{lag_plot.png}
    \caption{Лаговые диаграммы количественных признаков}
    \label{fig:lag}
\end{figure}

На Lag Plot наблюдается определённая структура точек. Соседние значения часто похожи друг на друга. Однако такая зависимость не является временной автокорреляцией, поскольку данные не представляют собой временной ряд. Структура возникает из-за порядка строк в датасете: наблюдения одного типа вина расположены последовательно.

\subsubsection*{Гистограммы}

Гистограммы используются для анализа формы распределения признаков, выявления асимметрии, нескольких пиков и отдельных областей концентрации значений.

\begin{Code}
fig, axes = plt.subplots(4, 3, figsize=(15, 14))
axes = axes.flatten()
for i, col in enumerate(features):
	ax = axes[i]
	ax.hist(df[col], bins=20, edgecolor='black', alpha=0.8)
	ax.set_title(f'Гистограмма: {col}')
	ax.set_xlabel(col)
	ax.set_ylabel('Частота')
plt.tight_layout()
plt.show()
\end{Code}

\begin{figure}[H]
    \centering
    \includegraphics[width=0.9\textwidth]{histograms.png}
    \caption{Гистограммы количественных признаков}
    \label{fig:histograms}
\end{figure}

Гистограммы показывают, что распределения большинства признаков являются асимметричными и имеют несколько локальных пиков. Это связано с тем, что в общей выборке объединены два типа вина, имеющие различные физико-химические свойства. Следовательно, для общей выборки предположение о фиксированном распределении нарушается.

\subsubsection*{Normal Probability Plot}

Normal Probability Plot, или Q–Q-график, используется для проверки соответствия распределения признака нормальному закону.

\begin{Code}
fig, axes = plt.subplots(4, 3, figsize=(15, 14))
axes = axes.flatten()
for i, col in enumerate(features):
	ax = axes[i]
	stats.probplot(df[col], dist='norm', plot=ax)
	ax.set_title(f'Q-Q Plot: {col}')
plt.tight_layout()
plt.show()
\end{Code}

\begin{figure}[H]
    \centering
    \includegraphics[width=0.9\textwidth]{normal_probability_plot.png}
    \caption{Normal Probability Plot количественных признаков}
    \label{fig:normal-probability}
\end{figure}

Для большинства признаков точки значительно отклоняются от прямой. Это связано с наличием двух групп наблюдений, соответствующих разным типам вина. Следовательно, общее распределение этих признаков нельзя считать нормальным.


\subsubsection*{Вывод по части 2}

В результате проверки базовых предположений было установлено, что наблюдения датасета \texttt{Wine Quality} расположены неслучайно. Это связано с тем, что объекты одного типа вина идут последовательно. Гистограммы и \texttt{Normal Probability Plot} показывают, что объединённые распределения признаков не являются нормальными. Это объясняется смешением двух различных групп объектов. Поэтому предположения о случайности, постоянстве среднего уровня, фиксированной вариации и нормальности общей выборки не выполняются полностью
 
\subsection*{\textbf{Часть 3. Описательная статистика}}
\addcontentsline{toc}{subsection}{Часть 3. Описательная статистика}
 
На данном этапе для количественных признаков датасета Wine Quality были рассчитаны меры центральной тенденции и меры разброса. Анализ проводился для 11 признаков.
 
\subsubsection*{Меры центральной тенденции}
 
Для описания центрального положения распределений были рассчитаны среднее арифметическое, медиана и мода.
 
\begin{Code}
mean_values = df[features].mean()
median_values = df[features].median()
mode_values = df[features].mode().iloc[0]

central_tendency = pd.DataFrame({
	'Среднее': mean_values,
	'Медиана': median_values,
	'Мода': mode_values
})
display(central_tendency.round(3))
\end{Code}

\begin{table}[ht]
	\centering
	\caption{Средние значения количественных признаков}
	\label{tab:means-wine}
	\begin{tabularx}{0.65\textwidth}{LCCC}
		\hline
		Признак & Среднее значение & Медиана & Мода \\
		\hline
		\texttt{fixed acidity} & 7{,}215 & 7{,}000 & 6{,}800\\
		\texttt{volatile acidity}  & 0{,}340 & 0{,}290 & 0{,}280\\
		\texttt{citric acid} & 0{,}319 & 0{,}310 & 0{,}300\\
		\texttt{residual sugar}  & 5{,}443 & 3{,}000 & 2{,}000\\
		\texttt{chlorides}  & 0{,}056 & 0{,}047 & 0{,}044\\
		\texttt{free sulfur dioxide}  & 30{,}525 & 29{,}000 & 29{,}000\\
		\texttt{total sulfur dioxide}  & 115{,}745 & 118{,}000 & 111{,}000\\
		\texttt{desity}  & 0{,}995 & 0{,}995 & 0{,}997\\
		\texttt{pH}  & 3{,}219 & 3{,}210 & 3{,}160\\
		\texttt{sulphates}  & 0{,}531 & 0{,}510 & 0{,}500\\
		\texttt{alcohol}  & 10{,}492 & 10{,}300 & 9{,}500\\
		\hline
	\end{tabularx}
\end{table}
\newpage
Среднее арифметическое характеризует общий уровень значений признака.
Медиана делит упорядоченный ряд наблюдений на две равные части, а мода
показывает наиболее часто встречающееся значение.
 
\subsubsection*{Меры разброса}
 
Для оценки изменчивости признаков были рассчитаны дисперсия, стандартное отклонение, межквартильный размах, медианное абсолютное отклонение, размах и коэффициент вариации.
 
\begin{Code}
variance = df[features].var(ddof=1)
std = df[features].std(ddof=1)
iqr = (df[features].quantile(0.75) - df[features].quantile(0.25))
mad = (df[features].sub(df[features].median()).abs().median())
data_range = df[features].max() - df[features].min()
cv = std / df[features].mean() * 100

dispersion = pd.DataFrame({
	'Дисперсия': variance,
	'Стандартное отклонение': std,
	'IQR': iqr,
	'MAD': mad,
	'Размах': data_range,
	'Коэффициент вариации, %': cv
})
display(dispersion.round(3))
\end{Code}
 
Основная аналитическая таблица оформлена в соответствии с требованиями
из методических указаний.
 

 \begin{table}[H]
 	\centering
 	\caption{Аналитические показатели количественных признаков}
 	\label{tab:analytical-statistics}
 	\scriptsize
 	
 	\begin{tabularx}{\textwidth}{LCCCCCC}
 		\hline
 		Переменная &
 		Среднее &
 		Дисперсия &
 		Ст. откл. &
 		IQR &
 		MAD &
 		Размах \\
 		\hline
 		
 		\texttt{fixed acidity} &
 		7{,}215 &
 		1{,}681 &
 		1{,}296 &
 		1{,}300 &
 		0{,}600 &
 		12{,}100 \\
 		
 		\texttt{volatile acidity} &
 		0{,}340 &
 		0{,}027 &
 		0{,}165 &
 		0{,}170 &
 		0{,}080 &
 		1{,}500 \\
 		
 		\texttt{citric acid} &
 		0{,}319 &
 		0{,}021 &
 		0{,}145 &
 		0{,}140 &
 		0{,}070 &
 		1{,}660 \\
 		
 		\texttt{residual sugar} &
 		5{,}443 &
 		22{,}637 &
 		4{,}758 &
 		6{,}300 &
 		1{,}700 &
 		65{,}200 \\
 		
 		\texttt{chlorides} &
 		0{,}056 &
 		0{,}001 &
 		0{,}035 &
 		0{,}027 &
 		0{,}011 &
 		0{,}602 \\
 		
 		\texttt{free sulfur dioxide} &
 		30{,}525 &
 		315{,}041 &
 		17{,}749 &
 		24{,}000 &
 		12{,}000 &
 		288{,}000 \\
 		
 		\texttt{total sulfur dioxide} &
 		115{,}745 &
 		3194{,}720 &
 		56{,}522 &
 		79{,}000 &
 		39{,}000 &
 		434{,}000 \\
 		
 		\texttt{density} &
 		0{,}995 &
 		0{,}000 &
 		0{,}003 &
 		0{,}005 &
 		0{,}002 &
 		0{,}052 \\
 		
 		\texttt{pH} &
 		3{,}219 &
 		0{,}026 &
 		0{,}161 &
 		0{,}210 &
 		0{,}110 &
 		1{,}290 \\
 		
 		\texttt{sulphates} &
 		0{,}531 &
 		0{,}022 &
 		0{,}149 &
 		0{,}170 &
 		0{,}080 &
 		1{,}780 \\
 		
 		\texttt{alcohol} &
 		10{,}492 &
 		1{,}423 &
 		1{,}193 &
 		1{,}800 &
 		0{,}900 &
 		6{,}900 \\
 		
 		\hline
 	\end{tabularx}
 \end{table}
 
\subsubsection*{Сравнение среднего и медианы}
 
Для большинства признаков среднее значение и медиана различаются. Это указывает на асимметричность их распределений.
 
\subsubsection*{Интерпретация показателей разброса}
 

Наибольший коэффициент вариации наблюдается у признаков \texttt{residual sugar}, \texttt{chlorides} и \texttt{total sulfur dioxide}. Это свидетельствует о наибольшей относительной изменчивости данных признаков. Для признака \texttt{density} коэффициент вариации наименьший. При этом показатели разброса рассчитаны для объединённой выборки. Поэтому их следует интерпретировать с учётом того, что датасет содержит два различных типа вина.

\subsubsection*{Вывод по части 3}
 
В результате расчёта описательных статистик было установлено, что для большинства признаков среднее и медиана различаются, что указывает на асимметричную форму их распределений. Наибольшая относительная вариативность наблюдается у признаков, связанных с сахаром и хлоридами. Полученные результаты подтверждают необходимость дальнейшего сравнительного анализа распределений отдельно по типам вина.
 
\subsection*{\textbf{Часть 4. Визуализация распределений}}
\addcontentsline{toc}{subsection}{Часть 4. Визуализация распределений}
 
На данном этапе были исследованы распределения количественных признаков
датасета Wine Quality с помощью гистограмм, Box Plot, эмпирической функции
распределения и ядерной оценки плотности. Также были рассчитаны
процентили и выполнена проверка потенциальных выбросов по правилу
$1{,}5 \cdot IQR$.
 
\subsubsection*{Гистограммы с различным количеством бинов}
 
Для исследования влияния количества интервалов на форму гистограммы были использованы значения количества бинов, равные 10, 20 и 30.
 
\begin{Code}
bins_list = [10, 20, 30]

fig, axes = plt.subplots(len(features), len(bins_list), figsize=(15, 32))

for i, col in enumerate(features):
for j, b in enumerate(bins_list):
	ax = axes[i, j]
	ax.hist(df[col], bins=b, edgecolor='black', alpha=0.8)
	ax.set_title(f'{col}, bins = {b}')
	ax.set_xlabel(col)
	ax.set_ylabel('Частота')

plt.tight_layout()
plt.show()
\end{Code}
 
\begin{figure}[H]
    \centering
    \includegraphics[width=0.85\textwidth]{Part1.png}
    \label{fig:histograms-bins1}
\end{figure}
\begin{figure}[H]
	\centering
	\includegraphics[width=0.85\textwidth]{Part2.png}
	\includegraphics[width=0.85\textwidth]{Part3.png}
	\caption{Гистограммы количественных признаков при различном количестве бинов}
	\label{fig:histograms-bins2}
\end{figure}

При небольшом количестве бинов гистограмма получается более сглаженной, однако отдельные особенности распределения могут быть скрыты. При увеличении количества бинов становятся заметны дополнительные пики и промежутки между областями концентрации наблюдений.
 
\subsubsection*{Box Plot и обнаружение выбросов}
 
Box Plot позволяет одновременно сравнить медиану, межквартильный размах и потенциальные выбросы количественных признаков.
 
\begin{Code}
fig, axes = plt.subplots(4, 3, figsize=(14, 14))
axes = axes.flatten()

for i, col in enumerate(features):
	ax = axes[i]
	sns.boxplot(y=df[col], ax=ax, color='steelblue')
	ax.set_title(f'Box Plot: {col}')
	ax.set_ylabel('Значение')
	ax.set_xlabel(col)
	ax.grid(axis='y', alpha=0.3)

# если признаков меньше, чем ячеек — убираем пустые
for j in range(len(features), len(axes)):
	fig.delaxes(axes[j])

plt.tight_layout()
plt.show()
\end{Code}
 
\begin{figure}[H]
    \centering
    \includegraphics[width=0.8\textwidth]{BoxPlot1.png}
    \caption{Box Plot количественных признаков}
    \label{fig:BoxPlot1}
\end{figure}
 
Для обнаружения потенциальных выбросов использовалось правило
$1{,}5 \cdot IQR$. Нижняя и верхняя границы определялись по формулам:
 
\[
L = Q_1 - 1{,}5 \cdot IQR,
\]
 
\[
U = Q_3 + 1{,}5 \cdot IQR,
\]
 
где $Q_1$ --- нижний квартиль, $Q_3$ --- верхний квартиль,
а $IQR = Q_3-Q_1$ --- межквартильный размах.
 
\begin{Code}
Вoutliers = {}
outliers = {}
for feature in features:
	q1 = df[feature].quantile(0.25)
	q3 = df[feature].quantile(0.75)
	iqr = q3 - q1
	lower_bound = q1 - 1.5 * iqr
	upper_bound = q3 + 1.5 * iqr
	mask = (df[feature] < lower_bound) | (df[feature] > upper_bound)
	outliers[feature] = df.loc[mask, feature]
	print(f'\nПризнак: {feature}')
	print(f'Q1 = {q1:.3f}')
	print(f'Q3 = {q3:.3f}')
	print(f'IQR = {iqr:.3f}')
	print(f'Нижняя граница = {lower_bound:.3f}')
	print(f'Верхняя граница = {upper_bound:.3f}')
	print(f'Количество выбросов: {mask.sum()}')
\end{Code}

\begin{table}[H]
	\centering
	\caption{Границы выбросов количественных признаков}
	\label{tab:outliers}
	\scriptsize
	
	\begin{tabularx}{\textwidth}{LCCCCCC}
		\hline
		Признак &
		$Q_1$ &
		$Q_3$ &
		IQR &
		\makecell{Нижняя\\граница} &
		\makecell{Верхняя\\граница} &
		\makecell{Количество\\выбросов} \\
		\hline
		
		\texttt{fixed acidity} &
		6{,}400 &
		7{,}700 &
		1{,}300 &
		4{,}450 &
		9{,}650 &
		357 \\
		
		\texttt{volatile acidity} &
		0{,}230 &
		0{,}400 &
		0{,}170 &
		$-0{,}025$ &
		0{,}655 &
		377 \\
		
		\texttt{citric acid} &
		0{,}250 &
		0{,}390 &
		0{,}140 &
		0{,}040 &
		0{,}600 &
		509 \\
		
		\texttt{residual sugar} &
		1{,}800 &
		8{,}100 &
		6{,}300 &
		$-7{,}650$ &
		17{,}550 &
		118 \\
		
		\texttt{chlorides} &
		0{,}038 &
		0{,}065 &
		0{,}027 &
		$-0{,}003$ &
		0{,}106 &
		286 \\
		
		\texttt{free sulfur dioxide} &
		17{,}000 &
		41{,}000 &
		24{,}000 &
		$-19{,}000$ &
		77{,}000 &
		62 \\
		
		\texttt{total sulfur dioxide} &
		77{,}000 &
		156{,}000 &
		79{,}000 &
		$-41{,}500$ &
		274{,}500 &
		10 \\
		
		\texttt{density} &
		0{,}992 &
		0{,}997 &
		0{,}005 &
		0{,}985 &
		1{,}004 &
		3 \\
		
		\texttt{pH} &
		3{,}110 &
		3{,}320 &
		0{,}210 &
		2{,}795 &
		3{,}635 &
		73 \\
		
		\texttt{sulphates} &
		0{,}430 &
		0{,}600 &
		0{,}170 &
		0{,}175 &
		0{,}855 &
		191 \\
		
		\texttt{alcohol} &
		9{,}500 &
		11{,}300 &
		1{,}800 &
		6{,}800 &
		14{,}000 &
		3 \\
		
		\hline
	\end{tabularx}
\end{table}
Потенциальные выбросы были обнаружены у большинства признаков. Обнаруженные значения не следует автоматически удалять. Box Plot выделяет потенциальные выбросы только с точки зрения статистического правила, но они могут быть корректными наблюдениями для отдельных типов вина.

\subsubsection*{Эмпирическая функция распределения и процентили}
 
Эмпирическая функция распределения, или CDF, позволяет определить долю наблюдений ниже заданного значения и сравнить распределения между собой.
 
\begin{Code}
fig, axes = plt.subplots(4, 3, figsize=(15, 14))
axes = axes.flatten()
for i, col in enumerate(features):
ax = axes[i]
sorted_data = np.sort(df[col])
yvals = np.arange(len(sorted_data)) / float(len(sorted_data) - 1)
	ax.step(sorted_data, yvals, where='post')
	ax.set_title(f'Эмпирическая CDF: {col}')
	ax.set_xlabel(col)
	ax.set_ylabel('F(x)')
	ax.grid(alpha=0.3)
plt.tight_layout()
plt.show()
\end{Code}

\begin{figure}[H]
    \centering
    \includegraphics[width=0.85\textwidth]{CDF1.png}
    \caption{Эмпирические функции распределения количественных признаков}
    \label{fig:CDF1}
\end{figure}
 
 
\begin{Code}
	percentiles = [10, 25, 50, 75, 90]
	percentile_table = df[features].quantile(np.array(percentiles) / 100)
	percentile_table.index = [f'{p}-й процентиль' for p in percentiles]
	display(percentile_table.round(3))
\end{Code}

\begin{table}[H]
	\centering
	\caption{Процентили количественных признаков}
	\label{tab:percentiles}
	\scriptsize
	
	\begin{tabularx}{\textwidth}{LCCCCCCCCCCC}
		\hline
		Процентиль &
		\texttt{fixed acidity} &
		\texttt{volatile acidity} &
		\texttt{citric acid} &
		\texttt{residual sugar} &
		\texttt{chlorides} &
		\texttt{free sulfur dioxide} &
		\texttt{total sulfur dioxide} &
		\texttt{density} &
		\texttt{pH} &
		\texttt{sulphates} &
		\texttt{alcohol} \\
		\hline
		
		10-й &
		6{,}0 & 0{,}18 & 0{,}14 & 1{,}3 & 0{,}031 & 9{,}0 & 30{,}0 & 0{,}991 & 3{,}02 & 0{,}37 & 9{,}1 \\
		
		25-й &
		6{,}4 & 0{,}23 & 0{,}25 & 1{,}8 & 0{,}038 & 17{,}0 & 77{,}0 & 0{,}992 & 3{,}11 & 0{,}43 & 9{,}5 \\
		
		50-й &
		7{,}0 & 0{,}29 & 0{,}31 & 3{,}0 & 0{,}047 & 29{,}0 & 118{,}0 & 0{,}995 & 3{,}21 & 0{,}51 & 10{,}3 \\
		
		75-й &
		7{,}7 & 0{,}40 & 0{,}39 & 8{,}1 & 0{,}065 & 41{,}0 & 156{,}0 & 0{,}997 & 3{,}32 & 0{,}60 & 11{,}3 \\
		
		90-й &
		8{,}8 & 0{,}59 & 0{,}49 & 13{,}0 & 0{,}086 & 54{,}0 & 188{,}0 & 0{,}998 & 3{,}42 & 0{,}72 & 12{,}3 \\
		
		\hline
	\end{tabularx}
\end{table}


\subsubsection*{KDE и сравнение с гистограммой}
 
Для сглаженной оценки плотности распределений использовалась ядерная оценка плотности KDE..
 
\begin{Code}
fig, axes = plt.subplots(4, 3, figsize=(15, 14))
axes = axes.flatten()
for i, feature in enumerate(features):
	ax = axes[i]
	sns.histplot(data=df, x=feature, bins=20, stat='density', kde=True, ax=ax)
	ax.set_title(f'Гистограмма и KDE: {feature}')
	ax.set_xlabel(feature)
	ax.set_ylabel('Плотность')
plt.tight_layout()
plt.show()
\end{Code}
 
\begin{figure}[H]
    \centering
    \includegraphics[width=0.85\textwidth]{KDE1.png}
    \caption{Сравнение гистограмм и KDE}
    \label{fig:KDE1}
\end{figure}
 
KDE подтверждает результаты, полученные с помощью гистограмм. Для многих признаков наблюдается несколько пиков, соответствующих различным группам вина.
 
\subsubsection*{Вывод по части 4}
 
В результате визуального анализа было установлено, что изменение количества бинов влияет на детализацию гистограммы. Box Plot выявил потенциальные выбросы у большинства признаков. CDF и KDE показали существенные различия между распределениями. Следовательно, для более точной интерпретации распределения следует анализировать не только по общей выборке, но и отдельно по каждому типу вина.
 
\subsection*{\textbf{Часть 5. Анализ взаимосвязей}}
\addcontentsline{toc}{subsection}{Часть 5. Анализ взаимосвязей}
 
На данном этапе были исследованы взаимосвязи между количественными признаками датасета Wine Quality. Для этого использовались диаграммы рассеяния, линии линейной регрессии, ковариация, коэффициент корреляции Пирсона, матрица диаграмм рассеяния, а также проверялось влияние потенциальных выбросов на корреляционные связи.
 
\subsubsection*{Диаграммы рассеяния с линией регрессии}
 
Для анализа парных взаимосвязей между признаками были построены диаграммы рассеяния с линиями линейной регрессии.
 
\begin{Code}
pairs = [
	('fixed acidity', 'volatile acidity'),
	('fixed acidity', 'citric acid'),
	('fixed acidity', 'residual sugar'),
	('volatile acidity', 'citric acid'),
	('alcohol', 'quality'),
	('density', 'alcohol')
]

fig, axes = plt.subplots(2, 3, figsize=(17, 10))
for ax, (x, y) in zip(axes.flat, pairs):
	sns.regplot(data=df, x=x, y=y, ax=ax, scatter_kws={'alpha': 0.65}, line_kws={'color': 'red'})
	ax.set_title(f'{x} и {y}')
plt.tight_layout()
plt.show()
\end{Code}
 
\begin{figure}[H]
    \centering
    \includegraphics[width=0.95\textwidth]{Scatter.png}
    \caption{Диаграммы рассеяния с линиями линейной регрессии}
    \label{fig:Scatter}
\end{figure}
 
\subsubsection*{Ковариация}
 
Ковариация характеризует направление совместного изменения двух
признаков. Положительное значение ковариации указывает на
увеличение одного признака при увеличении другого, а отрицательное
значение --- на противоположное направление изменения.
 
\begin{Code}
covariance_matrix = df[features].cov()
display(covariance_matrix.round(3))
\end{Code}

\begin{figure}[H]
	\centering
	\includegraphics[width=0.95\textwidth]{Kov.png}
	\caption{Ковариационная матрица количественных признаков}
	\label{fig:Kov}
\end{figure}

\subsubsection*{Корреляция Пирсона}
 
Коэффициент корреляции Пирсона показывает силу и направление линейной связи между двумя количественными признаками.
Для расчёта корреляционной матрицы использовался следующий код:
 
\begin{Code}
pearson_matrix = df[features].corr(method='pearson')
display(pearson_matrix.round(3))

plt.figure(figsize=(10, 8))
sns.heatmap(pearson_matrix, annot=True, fmt='.2f', cmap='coolwarm', vmin=-1, vmax=1, square=True)
plt.title('Корреляция Пирсона')
plt.show()
\end{Code}
 
\begin{figure}[H]
    \centering
    \includegraphics[width=0.85\textwidth]{Heat.png}
    \caption{Тепловая карта корреляций Пирсона}
    \label{fig:Heat}
\end{figure}
 
\subsubsection*{Матрица диаграмм рассеяния}
 
Для одновременного анализа всех пар признаков была построена матрица диаграмм рассеяния с разделением наблюдений по типам вина.
 
\begin{Code}
sns.pairplot(df, vars=features, hue='type', diag_kind='hist')
plt.show()
\end{Code}
 
\begin{figure}[H]
    \centering
    \includegraphics[width=0.95\textwidth]{pairplot.png}
    \caption{Матрица диаграмм рассеяния количественных признаков}
    \label{fig:pairplot}
\end{figure}
 
\subsubsection*{Проверка влияния потенциальных выбросов}
 
В части 4 потенциальные выбросы были обнаружены у большинства признаков. Для проверки их влияния на корреляции были рассчитаны корреляционные матрицы до и после исключения таких наблюдений.
 
\begin{Code}
outlier_mask = pd.Series(False, index=df.index)
for feature in features:
q1 = df[feature].quantile(0.25)
q3 = df[feature].quantile(0.75)
iqr = q3 - q1
lower_bound = q1 - 1.5 * iqr
upper_bound = q3 + 1.5 * iqr
feature_outliers = (df[feature] < lower_bound) | (df[feature] > upper_bound)
outlier_mask = outlier_mask | feature_outliers

print('Количество наблюдений с потенциальными выбросами:', outlier_mask.sum())

corr_all = df[features].corr()
corr_without_outliers = df.loc[~outlier_mask, features].corr()

print('Корреляция для всей выборки:')
display(corr_all.round(3))
print('Корреляция без выбросов:')
display(corr_without_outliers.round(3))
\end{Code}
 
\begin{figure}[H]
	\centering
	\includegraphics[width=0.95\textwidth]{Kor1.png}
	\includegraphics[width=0.95\textwidth]{Kor2.png}
	\caption{Корреляция для всей выборки и без выбросов}
	\label{fig:Kor}
\end{figure}

\subsubsection*{Вывод по части 5}
 
В результате анализа взаимосвязей было установлено, что наиболее сильная положительная корреляция наблюдается между \texttt{free sulfur dioxide} и \texttt{total sulfur dioxide}. Также наблюдается положительная связь между \texttt{alcohol} и \texttt{quality}. Матрица диаграмм рассеяния показала, что признаки \texttt{alcohol} и \texttt{volatile acidity} лучше всего разделяют типы вина. Проверка влияния потенциальных выбросов показала, что их исключение незначительно изменяет значения коэффициентов корреляции.
 
\subsection*{\textbf{Часть 6. Сравнительный анализ}}
\addcontentsline{toc}{subsection}{Часть 6. Сравнительный анализ}
 
На данном этапе был проведён сравнительный анализ распределений количественного признака по группам, заданным категориальной переменной.
 
\subsubsection*{Выбор категориальной и количественной переменных}
 
В качестве категориальной переменной был выбран признак \texttt{type}, содержащий названия двух типов вина: \texttt{red} и \texttt{white}. В качестве количественной переменной был выбран признак \texttt{alcohol}, характеризующий содержание алкоголя в процентах.
 
\subsubsection*{Box Plot по группам}
 
Для сравнения распределений содержания алкоголя между типами вина были построены Box Plot по группам.
 
\begin{Code}
plt.figure(figsize=(10, 6))
sns.boxplot(data=df, x='type', y='alcohol')
plt.title('Сравнение содержания алкоголя по типам вина')
plt.xlabel('Тип вина')
plt.ylabel('Содержание алкоголя')
plt.grid(axis='y', alpha=0.3)
plt.show()
\end{Code}
 
\begin{figure}[H]
    \centering
    \includegraphics[width=0.85\textwidth]{BoxPl.png}
    \caption{Сравнение содержания алкоголя по типам вина}
    \label{fig:BoxPl}
\end{figure}
 
\subsubsection*{Эмпирическая функция распределения по группам}
 
Для дополнительного сравнения групп были построены эмпирические функции распределения CDF.
 
\begin{Code}
plt.figure(figsize=(10, 6))
for wine_type in df['type'].unique():
	values = np.sort(df.loc[df['type'] == wine_type, 'alcohol'])
	probabilities = np.arange(1, len(values) + 1) / len(values)
	plt.step(values, probabilities, where='post', label=wine_type)
plt.title('CDF для содержания алкоголя по типам вина')
plt.xlabel('Содержание алкоголя')
plt.ylabel('F(x)')
plt.ylim(0, 1.05)
plt.grid(alpha=0.3)
plt.legend()
plt.show()
\end{Code}
 
\begin{figure}[H]
    \centering
    \includegraphics[width=0.85\textwidth]{CDf3.png}
    \caption{CDF для содержания алкоголя по типам вина}
    \label{fig:CDf3}
\end{figure}
 
\subsubsection*{Сравнение медиан и разброса}
 
Для количественного сравнения групп были рассчитаны среднее значение,
медиана, стандартное отклонение, минимальное и максимальное значения,
первый и третий квартили, а также межквартильный размах.
 
\begin{Code}
group_statistics = df.groupby('type')['alcohol'].agg(
	Среднее='mean',
	Медиана='median',
	Стандартное_отклонение='std',
	Минимум='min',
	Максимум='max',
	Q1=lambda x: x.quantile(0.25),
	Q3=lambda x: x.quantile(0.75)
)
group_statistics['IQR'] = group_statistics['Q3'] - group_statistics['Q1']
display(group_statistics.round(3))
\end{Code}
 
\begin{table}[H]
	\centering
	\caption{Сравнительные статистики по типу вина}
	\label{tab:by-type}
	\small
	
	\begin{tabular}{lcccccccc}
		\hline
		\textbf{type} &
		\textbf{Среднее} &
		\textbf{Медиана} &
		\makecell{\textbf{Стандартное}\\\textbf{отклонение}} &
		\textbf{Минимум} &
		\textbf{Максимум} &
		\textbf{Q1} &
		\textbf{Q3} &
		\textbf{IQR} \\
		\hline
		
		red   & 10{,}423 & 10{,}2 & 1{,}066 & 8{,}4 & 14{,}9 & 9{,}5 & 11{,}1 & 1{,}6 \\
		white & 10{,}514 & 10{,}4 & 1{,}231 & 8{,}0 & 14{,}2 & 9{,}5 & 11{,}4 & 1{,}9 \\
		\hline
	\end{tabular}
\end{table}
 
\subsubsection*{Интерпретация различий между группами}
 
Результаты Box Plot, CDF и численного сравнения показывают, что типы вина имеют существенно различающиеся распределения признака alcohol. Белое вино характеризуется несколько более высоким средним и медианным содержанием алкоголя, а также большим разбросом значений. Красное вино имеет более низкие и более сгруппированные значения.
 
\subsubsection*{Вывод по части 6}
 
В качестве категориальной переменной был выбран признак \texttt{type}, а в качестве количественной переменной — \texttt{alcohol}. Сравнительный анализ показал, что медианные значения содержания алкоголя существенно различаются между типами вина. Следовательно, признак \texttt{alcohol} позволяет хорошо различать типы вина и может рассматриваться как информативный признак для дальнейшего моделирования или классификации.











